% Terjemahan Bahasa Indonesia dari prelude.tex baris 207-495.
% Sumber: Methods of Algebra, Volume 2, commit
% 9a5803ff77dd3257484cb177f851a73770a59dd3 (CC BY 4.0).
% stable-unit-id: o014.aljabr2.prelude.conventions

% segment-id: o014.li2.prelude.conventions.h001
\section*{Konvensi}
\addcontentsline{toc}{section}{Konvensi}
\hypertarget{o014.aljabr2.prelude.conventions}{ }

% segment-id: o014.li2.prelude.conventions.p001
Konvensi notasi buku ini sama dengan Jilid I. Walaupun agak panjang, konvensi
tersebut dirangkum kembali di sini.

% segment-id: o014.li2.prelude.conventions.h002
\subsection*{Konvensi Dasar}

% segment-id: o014.li2.prelude.conventions.p002
Buku ini menggunakan simbol logika standar seperti $\forall$, $\exists$,
$\implies$, dan sebagainya. Simbol $\exists !$ berarti ``ada tepat satu'',
sedangkan $\iff$ berarti ``jika dan hanya jika''. Akhir suatu bukti ditandai
dengan $\Box$. Notasi $A := B$ berarti ``$A$ didefinisikan sebagai $B$''.
Apabila definisi suatu objek matematika tidak bergantung pada pilihan data
bantu, objek itu disebut terdefinisi dengan baik.
\index{terdefinisi dengan baik (well-defined)}

% segment-id: o014.li2.prelude.conventions.p003
Panah $\rightiso$ menyatakan isomorfisme antarobjek; kadang-kadang isomorfisme
itu ditulis tanpa arah sebagai $\simeq$. Notasi $\xrightarrow{1:1}$ atau
$\stackrel{1:1}{\longleftrightarrow}$ menyatakan bijeksi antardua himpunan,
yakni korespondensi satu-satu. Notasi berbentuk $F(\cdot)$ dimaksudkan untuk
menonjolkan peubah pada pemetaan atau funktor $F$.

% segment-id: o014.li2.prelude.conventions.p004
Notasi $n \gg m$ berarti bahwa $n$ cukup jauh lebih besar daripada $m$.
Sebagai contoh, $n \gg 0$ berarti bahwa $n$ cukup besar, sedangkan $n \ll 0$
berarti bahwa $-n$ cukup besar.

% segment-id: o014.li2.prelude.conventions.p005
Buku ini banyak menggunakan bahasa diagram komutatif. Kasus yang paling
sederhana ditunjukkan oleh diagram berikut.

% segment-id: o014.li2.prelude.conventions.d001
\[\begin{tikzcd}
	A \arrow[r, "f"] \arrow[d, "h"'] & B \arrow[ld, "g"] \\
	C &
\end{tikzcd} \; \text{komutatif} \stackrel{\text{definisi}}{\iff} g \circ f = h. \]

% segment-id: o014.li2.prelude.conventions.p006
Dalam diagram ini, $f$, $g$, dan $h$ dapat berupa pemetaan, homomorfisme, atau morfisme
dalam suatu kategori umum. Diagram komutatif yang lebih rumit dipahami dengan
cara yang serupa.

% segment-id: o014.li2.prelude.conventions.h003
\subsection*{Himpunan dan Struktur Urutan}

% segment-id: o014.li2.prelude.conventions.p007
Buku ini menggunakan teori himpunan aksiomatik Zermelo--Fraenkel dan menerima
aksioma pilihan.

% segment-id: o014.li2.prelude.conventions.p008
Himpunan kosong ditulis $\emptyset$. Notasi $A \subset B$ berarti bahwa $A$
merupakan himpunan bagian dari $B$, dengan kesamaan diperbolehkan; inklusi
sejati ditulis $A \subsetneq B$. Selisih himpunan ditulis
$A \smallsetminus B := \left\{a \in A: a \notin B \right\}$, sedangkan gabungan
saling lepas ditulis $A \sqcup B$. Untuk suatu pemetaan $f: A \to B$, notasi
$a \mapsto b$ berarti bahwa $a \in A$ dipetakan oleh $f$ ke $b \in B$, dan
$f|_{A'}$ menyatakan pembatasan $f$ pada himpunan bagian $A' \subset A$.
Pemetaan identitas pada $A$ ditulis $\identity = \identity_A$. Banyaknya unsur,
kardinalitas, atau kardinal himpunan $A$ ditulis $|A|$.

% segment-id: o014.li2.prelude.conventions.p009
Berdasarkan \cite[Definisi~1.2.1]{Li1}, sebuah \emph{himpunan terurut parsial}
\index{himpunan terurut parsial (partially ordered set)} ialah himpunan $P$
yang dilengkapi relasi biner $\leq$ yang refleksif, transitif, dan antisimetris.
Jika antisimetrisitas tidak disyaratkan, diperoleh \emph{himpunan
praterurut}. Notasi $x < y$ berarti $x \lneq y$. Himpunan terurut parsial $(P, \leq)$
sering disingkat menjadi $P$.

% segment-id: o014.li2.prelude.conventions.l001
\begin{itemize}
	% segment-id: o014.li2.prelude.conventions.i001
	\item Jika $m \in P$ memenuhi $\forall x \in P, \; x \geq m \iff x = m$,
	maka $m$ disebut unsur maksimal dalam $P$. Konsep unsur minimal didefinisikan
	dengan mengganti $\geq$ oleh $\leq$. Pada umumnya unsur-unsur ini tidak
	tunggal.
	% segment-id: o014.li2.prelude.conventions.i002
	\item Misalkan $S$ himpunan bagian dari $P$. Jika $b \in P$ memenuhi
	$\forall x \in S, \; x \leq b$, maka unsur itu disebut batas atas $S$ dalam $P$.
	Batas bawah $S$ didefinisikan secara serupa.
	% segment-id: o014.li2.prelude.conventions.i003
	\item Jika $b \in P$ merupakan batas atas $S$ dan setiap batas atas lain
	$b'$ memenuhi $b' \geq b$, maka unsur itu disebut supremum $S$, ditulis $\sup S$.
	Supremum, jika ada, bersifat tunggal. Infimum $S$, ditulis $\inf S$,
	didefinisikan secara serupa.
\end{itemize}

% segment-id: o014.li2.prelude.conventions.p010
Suatu pemetaan $f: P \to Q$ antarhimpunan terurut parsial disebut
mempertahankan urutan jika
$a \leq b \implies f(a) \leq f(b)$. Jika $f: P \to Q$ bijektif serta $f$ dan
$f^{-1}$ sama-sama mempertahankan urutan, maka $f$ disebut isomorfisme
antarhimpunan terurut parsial, atau bijeksi yang mempertahankan urutan.

% segment-id: o014.li2.prelude.conventions.p011
Jika setiap dua unsur himpunan terurut parsial $P$ dapat dibandingkan, maka
himpunan itu disebut himpunan terurut total. Jika setiap himpunan bagian tak kosong dari
himpunan terurut total tak kosong $P$ mempunyai unsur terkecil, maka $P$
disebut \emph{himpunan terurut baik}
\index{himpunan terurut baik (well-ordered set)}. Suatu kelas khusus himpunan
terurut baik disebut \emph{ordinal}\index{ordinal (ordinal)}, dan setiap dua
ordinal dapat dibandingkan. Ringkasnya:

% segment-id: o014.li2.prelude.conventions.l002
\begin{itemize}
	% segment-id: o014.li2.prelude.conventions.i004
	\item untuk setiap ordinal $\alpha$, berlaku
	$\alpha = \left\{ \beta: \text{ordinal}, \; \beta < \alpha \right\}$;
	% segment-id: o014.li2.prelude.conventions.i005
	\item jika ordinal $\kappa$, sebagai himpunan, tidak ekuipoten dengan
	ordinal mana pun yang $< \kappa$, maka ordinal itu disebut \emph{kardinal};
	\index{kardinal (cardinal)}
	% segment-id: o014.li2.prelude.conventions.i006
	\item kardinalitas $|S|$ suatu himpunan $S$ didefinisikan sebagai ordinal
	terkecil dalam kelas ekuipotensinya. Di sini digunakan fakta bahwa setiap
	himpunan dapat diberi urutan baik (yang memerlukan aksioma pilihan) dan bahwa
	setiap himpunan terurut baik isomorfik dengan tepat satu ordinal;
	% segment-id: o014.li2.prelude.conventions.i007
	\item khususnya, $|\alpha| \leq \alpha$ untuk setiap ordinal $\alpha$.
\end{itemize}

% segment-id: o014.li2.prelude.conventions.p012
Inilah sudut pandang von Neumann; lihat \cite[\S~1.2]{Li1} atau \cite{Je03}.
Pendekatan ini sedikit lebih berliku daripada definisi lazim kardinal sebagai
kelas ekuipotensi, tetapi juga mempunyai kegunaan tersendiri.

% segment-id: o014.li2.prelude.conventions.p013
Ordinal hingga dapat diidentifikasi dengan bilangan bulat tak negatif: untuk
setiap $n \in \Z_{\geq 0}$ terdapat himpunan terurut total $\mathbf{n}$, yang
sebagai himpunan didefinisikan secara rekursif oleh
$\mathbf{0} := \emptyset$ dan
$\mathbf{n+1} := \{\mathbf{0}, \ldots, \mathbf{n}\}$; struktur urutannya jelas.
Ordinal tak hingga terkecil $\omega$ ialah himpunan terurut baik
$\mathbf{0} \leq \mathbf{1} \leq \cdots$, dan sekaligus merupakan kardinal
terhitung $\aleph_0$.
\index[sym1]{012@$\mathbf{0}, \mathbf{1}, \mathbf{2}, \ldots$}

% segment-id: o014.li2.prelude.conventions.p014
Kita memilih suatu \emph{semesta Grothendieck}
\index{semesta Grothendieck (Grothendieck universe)}
\index[sym1]{U@$\mathcal{U}$} $\mathcal{U}$ untuk membedakan ukuran himpunan.
Unsur-unsur $\mathcal{U}$ disebut himpunan-$\mathcal{U}$; himpunan yang
ekuipoten dengan suatu himpunan-$\mathcal{U}$ disebut himpunan
$\mathcal{U}$-kecil, atau cukup \emph{himpunan kecil}
\index{himpunan kecil (small set)}. Jika kardinal $\mu$, sebagai himpunan,
merupakan himpunan $\mathcal{U}$-kecil, maka kardinal itu disebut \emph{kardinal
kecil}\index{kardinal!kecil (small)}. Kita mengandaikan bahwa setiap himpunan
$X$ termasuk dalam suatu $\mathcal{U}$; lihat \cite[Asumsi~1.5.2]{Li1} dan
pembahasan terkait.

% segment-id: o014.li2.prelude.conventions.h004
\subsection*{Kategori}

% segment-id: o014.li2.prelude.conventions.p015
Konvensi tentang kategori sama dengan \cite[\S~1.5, Definisi~2.1.3]{Li1}:
semua objek dan semua morfisme suatu kategori $\mathcal{C}$ masing-masing
membentuk himpunan, yang ditulis $\Obj(\mathcal{C})$ dan $\Mor(\mathcal{C})$.
Himpunan morfisme di antara dua objek ditulis
$\Hom_{\mathcal{C}}(\cdot, \cdot)$, atau cukup $\Hom(\cdot, \cdot)$. Morfisme
identitas objek $X$ ditulis $\identity_X$. Monomorfisme ditandai dengan panah
$\hookrightarrow$, sedangkan epimorfisme dengan panah $\twoheadrightarrow$;
monomorfisme juga disebut pembenaman.
\index{kategori (category)}
\index{objek (object)}

% segment-id: o014.li2.prelude.conventions.l003
\begin{itemize}
	% segment-id: o014.li2.prelude.conventions.i008
	\item Untuk setiap morfisme $f: X \to Y$ dalam suatu kategori dan setiap
	objek $T$, morfisme $f$ menginduksi operasi tarikan balik $f^*$ dan dorongan
	maju $f_*$ pada himpunan $\Hom$:
	\begin{align*}
		f^*: \Hom(Y, T) & \to \Hom(X, T), & f_*: \Hom(T, X) & \to \Hom(T, Y) \\
		\beta & \mapsto \beta f & \alpha & \mapsto f\alpha .
	\end{align*}
	% segment-id: o014.li2.prelude.conventions.i009
	\item Subkategori ditulis dengan notasi himpunan bagian,
	$\mathcal{C}' \subset \mathcal{C}$. Jika
	$\Hom_{\mathcal{C}'}(X, Y) = \Hom_{\mathcal{C}}(X, Y)$, maka
	$\mathcal{C}'$ disebut subkategori penuh.

	% segment-id: o014.li2.prelude.conventions.i010
	\item Produk dua kategori $\mathcal{C}_1$ dan $\mathcal{C}_2$,
	$\mathcal{C}_1 \times \mathcal{C}_2$, mempunyai himpunan objek
	$\Obj(\mathcal{C}_1) \times \Obj(\mathcal{C}_2)$. Morfisme dari
	$(X_1, X_2)$ ke $(Y_1, Y_2)$ adalah unsur
	$\Hom_{\mathcal{C}_1}(X_1, Y_1) \times
	\Hom_{\mathcal{C}_2}(X_2, Y_2)$, dan komposisi morfismenya dilakukan per
	komponen. Dari sini terlihat cara mendefinisikan produk sebarang keluarga
	kategori $\prod_i \mathcal{C}_i$. Dengan cara yang sama didefinisikan
	$\mathcal{C}^n := \mathcal{C} \times \cdots \times \mathcal{C}$ (sebanyak
	$n$ faktor), atau lebih umum $\mathcal{C}^I$, dengan $I$ suatu himpunan.
	Lihat \cite[Definisi~2.3.1]{Li1}.
	\index{kategori produk (product category)}
\end{itemize}

% segment-id: o014.li2.prelude.conventions.p016
Sebagai contoh, setiap himpunan praterurut $(P, \leq)$ menentukan kategori
$\mathcal{P}$ dengan himpunan objek $P$ dan

% segment-id: o014.li2.prelude.conventions.d002
\[ \forall a,b \in P, \quad \Hom_{\mathcal{P}}(a, b) = \begin{cases}
	\text{himpunan satu titik} \; \{ \star \}, & a \leq b \\
	\emptyset, & \text{selain itu.}
\end{cases}\]

% segment-id: o014.li2.prelude.conventions.p017
Sebagai kasus khusus, $\mathbf{0}$ dapat diidentifikasi dengan kategori kosong,
tanpa objek dan tanpa morfisme, sedangkan $\mathbf{1}$ diidentifikasi dengan
kategori dengan tepat satu objek $\mathbf{0}$ dan satu morfisme
$\identity_{\mathbf{0}}$.

% segment-id: o014.li2.prelude.conventions.p018
Kategori yang memenuhi
$\Mor(\mathcal{C}) = \left\{ \identity_X : X \in \Obj(\mathcal{C}) \right\}$
disebut \emph{kategori diskret}. Kategori-kategori ini berkorespondensi satu-satu
dengan himpunan melalui $\mathcal{C} \mapsto \Obj(\mathcal{C})$.
\index{kategori!diskret (discrete)}

% segment-id: o014.li2.prelude.conventions.h005
\subsection*{Terminologi: Kategori Kecil dan Kategori Besar}

% segment-id: o014.li2.prelude.conventions.p019
Buku ini memilih suatu semesta Grothendieck $\mathcal{U}$. Misalkan
$\mathcal{C}$ suatu kategori. Jika $\mathrm{Mor}(\mathcal{C})$ merupakan
himpunan $\mathcal{U}$-kecil, maka kategori itu disebut kategori
$\mathcal{U}$-kecil. Jika hanya disyaratkan bahwa $\Hom_{\mathcal{C}}(X, Y)$
merupakan himpunan $\mathcal{U}$-kecil untuk semua
$X, Y \in \Obj(\mathcal{C})$, maka kategori itu disebut kategori
$\mathcal{U}$.
\footnote{Dalam banyak pustaka, objek-objek suatu kategori dipahami membentuk
kelas, sedangkan himpunan $\Hom$ di antara setiap dua objek merupakan himpunan.
Dengan demikian, ``kelas'' dalam pustaka tersebut berpadanan dengan himpunan
dalam buku ini, ``himpunan'' mereka berpadanan dengan himpunan-$\mathcal{U}$
dalam buku ini, dan ``kategori'' mereka pada dasarnya berpadanan dengan
kategori $\mathcal{U}$ dalam buku ini.}

% segment-id: o014.li2.prelude.conventions.p020
Sebagai contoh, semua himpunan-$\mathcal{U}$ beserta pemetaan di antaranya
membentuk kategori-$\mathcal{U}$ yang ditulis $\cate{Set}$; kategori ini bukan
kategori $\mathcal{U}$-kecil. Di sisi lain, tanpa memilih $\mathcal{U}$,
semua himpunan membentuk kelas sejati, bukan kategori dalam pengertian buku ini.
Serupa dengan itu, $\cate{Grp}$ (atau $\cate{Ab}$) dalam buku ini menyatakan
kategori semua grup (atau grup abelian) yang dibangun pada
himpunan-$\mathcal{U}$; semuanya merupakan kategori-$\mathcal{U}$.
\index[sym1]{Set@$\cate{Set}$}

% segment-id: o014.li2.prelude.conventions.p021
Untuk melihat bahwa ukuran benar-benar menimbulkan persoalan, perhatikan bahwa
funktor $\Hom$, yaitu
$\Hom_{\mathcal{C}}: \mathcal{C}^{\opp} \times \mathcal{C} \to \cate{Set}$,
hanya dapat didefinisikan apabila $\mathcal{C}$ merupakan kategori-$\mathcal{U}$.

% segment-id: o014.li2.prelude.conventions.p022
Kecuali dinyatakan lain, kategori yang dibahas dalam buku ini selalu berarti
kategori-$\mathcal{U}$, dan kategori $\mathcal{U}$-kecil cukup disebut
\emph{kategori kecil}. Kategori yang belum tentu merupakan kategori-$\mathcal{U}$
disebut \emph{kategori besar}.
\index{kategori!kecil, besar (small, big)}
\index{masalah ukuran (size issues)}

% segment-id: o014.li2.prelude.conventions.p023
Kategori besar sukar dihindari dalam banyak konstruksi, khususnya teori yang
berkaitan dengan kategori funktor atau lokalisasi kategoris. Setiap kali kategori besar
mungkin muncul dan benar-benar menimbulkan akibat substantif, hal itu akan
ditandai secara eksplisit.

% segment-id: o014.li2.prelude.conventions.h006
\subsection*{Funktor dan Kategori Funktor}

% segment-id: o014.li2.prelude.conventions.p024
Funktor identitas dari kategori $\mathcal{C}$ ke dirinya sendiri ditulis
$\identity_{\mathcal{C}}$. Suatu funktor $F: \mathcal{C} \to \mathcal{D}$
memetakan morfisme $f: X \to Y$ dalam $\mathcal{C}$ ke morfisme
$Ff: FX \to FY$ dalam $\mathcal{D}$. Jika pemetaan
$\Hom_{\mathcal{C}}(X, Y) \to \Hom_{\mathcal{D}}(FX, FY)$ bersifat injektif
(atau bijektif) untuk semua $X$ dan $Y$, maka $F$ disebut setia (atau penuh dan
setia). Jika setiap objek $\mathcal{D}$ isomorfik dengan suatu $FX$, maka $F$
disebut pada hakikatnya surjektif.
\index{funktor (functor)}

% segment-id: o014.li2.prelude.conventions.p025
Sebagai contoh, bagi kategori-kategori yang terkait dengan himpunan terurut
parsial, pemetaan yang mempertahankan urutan $P \to Q$ tidak lain ialah funktor
$\mathcal{P} \to \mathcal{Q}$.

\index{transformasi alami (natural transformation)}
\index{2-sel (2-cell)}
% segment-id: o014.li2.prelude.conventions.p026
Dalam sebagian pustaka, suatu morfisme $\varphi: F \to G$ di antara dua funktor
juga disebut transformasi alami, dan kadang-kadang ditulis dengan panah ganda
$\varphi: F \Rightarrow G$. Morfisme $\varphi$ terdiri atas keluarga morfisme
yang kompatibel $(\varphi_X: FX \to GX)_{X \in \Obj(\mathcal{C})}$, dan dapat
pula dinyatakan dengan diagram \emph{$2$-sel} yang diperkenalkan dalam
\cite[\S~2.2]{Li1}:

% segment-id: o014.li2.prelude.conventions.d003
\[\begin{tikzcd}
	\mathcal{C} \arrow[rr, bend left=45, "F", ""'{name=F}] \arrow[rr, bend right=45, "G"', ""{name=G}] & & \mathcal{D} \arrow[Rightarrow, from=F, to=G, "\varphi"].
\end{tikzcd}\]

% segment-id: o014.li2.prelude.conventions.p027
Ketika kelak membahas kategori-$\cate{Ab}$ (lihat
\S\sourcecrossref{sec:Ker-Coker}{1.3}) atau kategori $\Bbbk$-linear (lihat
\S\sourcecrossref{sec:linear-cat}{1.4}), funktor akan dikenai syarat tambahan
yang sesuai.

% segment-id: o014.li2.prelude.conventions.p028
Diberikan funktor-funktor
$\begin{tikzcd}
	\mathcal{B} \arrow[r, "R"] & \mathcal{C} \arrow[shift left, r, "F"] \arrow[shift right, r, "G"'] & \mathcal{D} \arrow[r, "L"] & \mathcal{E}
\end{tikzcd}$
dan morfisme $\varphi: F \to G$, secara alami diperoleh morfisme
$L\varphi: LF \to LG$ dan $\varphi R: FR \to GR$. Di sisi lain, morfisme
$\varphi: F \to G$ dan $\psi: G \to H$ dapat dikomposisikan menjadi
$\psi\varphi: F \to H$. Operasi-operasi komposisi ini dapat dinyatakan sebagai
komposisi $2$-sel, yakni dengan menempelkan diagram. Identitas segitiga yang
segera akan ditinjau kembali merupakan contoh dasar.

\index{kategori funktor (functor category)}
\index[sym1]{DC@$\mathcal{D}^{\mathcal{C}}$}
\index[sym1]{Fct@$\mathrm{Fct}$}
% segment-id: o014.li2.prelude.conventions.p029
Semua funktor dari kategori $\mathcal{C}$ ke kategori $\mathcal{D}$ membentuk
kategori funktor $\mathcal{D}^{\mathcal{C}}$, yang juga ditulis
$\mathrm{Fct}(\mathcal{C}, \mathcal{D})$. Perhatikan bahwa, kecuali
$\mathcal{C}$ kecil, $\mathcal{D}^{\mathcal{C}}$ sering kali merupakan
``kategori besar''. Sebagai kasus khusus kategori funktor, untuk
$n \in \Z_{\geq 0}$ dan himpunan terurut total terkait $\mathbf{n}$, kategori
$\mathcal{C}^{\mathbf{n}}$ dapat dijelaskan sebagai berikut.

% segment-id: o014.li2.prelude.conventions.l004
\begin{compactitem}
	% segment-id: o014.li2.prelude.conventions.i011
	\item menurut definisi, $\mathcal{C}^{\mathbf{0}} := \mathbf{1}$;
	% segment-id: o014.li2.prelude.conventions.i012
	\item terdapat tepat satu funktor $\mathcal{C} \to \mathbf{1}$, sehingga
	$\mathbf{1}^{\mathcal{C}} = \mathbf{1}$;
	% segment-id: o014.li2.prelude.conventions.i013
	\item menentukan funktor $\mathbf{1} \to \mathcal{C}$ setara dengan
	menentukan suatu objek $\mathcal{C}$, sehingga
	$\mathcal{C}^{\mathbf{1}} \simeq \mathcal{C}$;
	% segment-id: o014.li2.prelude.conventions.i014
	\item objek-objek $\mathcal{C}^{\mathbf{2}}$ ialah morfisme dalam
	$\mathcal{C}$, sedangkan morfisme-morfismenya ialah persegi komutatif di
	antara morfisme-morfisme tersebut.
\end{compactitem}

% segment-id: o014.li2.prelude.conventions.p030
Dengan demikian, $\mathbf{0}$ dapat disebut kategori awal,
$\mathbf{1} = [\bullet]$ kategori terminal, dan
$\mathbf{2} = [\bullet \to \bullet]$ dapat dibayangkan sebagai panah berjalan.

\index{kategori!lawan (opposite)}
\index[sym1]{Cop@$\mathcal{C}^{\opp}$}
% segment-id: o014.li2.prelude.conventions.p031
Kategori lawan $\mathcal{C}^{\opp}$ dari $\mathcal{C}$ memiliki himpunan objek
yang sama dengan $\mathcal{C}$. Setiap morfisme $f: X \to Y$ dapat dipandang
sebagai morfisme $Y \to X$ dalam $\mathcal{C}^{\opp}$ dan, untuk
membedakannya, ditulis $f^{\opp}$\index[sym1]{fop@$f^{\opp}$}. Di sisi lain,
$\mathrm{Fct}(\mathcal{C}, \mathcal{D})^{\opp}$ secara alami dapat
diidentifikasi dengan
$\mathrm{Fct}(\mathcal{C}^{\opp}, \mathcal{D}^{\opp})$. Peralihan dari
$\mathcal{C}$ ke $\mathcal{C}^{\opp}$ berarti membalik semua panah; mekanisme
ini disebut dualitas dalam teori kategori.

% segment-id: o014.li2.prelude.conventions.p032
Jika sepasang funktor $F: \mathcal{C} \leftrightarrows \mathcal{D}: G$
memenuhi $GF \simeq \identity_{\mathcal{C}}$ dan
$FG \simeq \identity_{\mathcal{D}}$, keduanya disebut saling kuasi-invers.
Funktor yang mempunyai kuasi-invers disebut ekuivalensi antarkategori, suatu
konsep dasar dalam teori kategori. Jika berlaku persamaan ketat
$GF = \identity_{\mathcal{C}}$ dan $FG = \identity_{\mathcal{D}}$, maka
keduanya disebut isomorfisme antarkategori yang saling invers.

% segment-id: o014.li2.prelude.conventions.h007
\subsection*{Struktur Aljabar}

% segment-id: o014.li2.prelude.conventions.p033
Unsur identitas (atau unsur satuan) suatu grup ditulis $1$ dalam notasi
multiplikatif, sedangkan unsur identitas grup abelian ditulis $0$ dalam notasi
aditif. Grup lawan dari $G$, yang urutan perkaliannya dibalik, ditulis
$G^{\opp}$. Jika $H$ subgrup dari $G$, indeks $H$ dalam $G$ ditulis $(G:H)$;
nilainya juga sama dengan banyaknya koset $|G/H|$ atau $|H \backslash G|$.
\index[sym1]{GH@$(G:H)$}

% segment-id: o014.li2.prelude.conventions.p034
Kecuali dinyatakan lain, gelanggang selalu berarti gelanggang asosiatif dengan
unsur satuan; demikian pula untuk aljabar atas gelanggang komutatif. Unsur
satuan perkalian gelanggang $R$ ditulis $1$ atau $1_R$, dan semua unsur
yang dapat dibalik membentuk grup perkalian $R^\times$. Karakteristik lapangan atau
daerah integral $R$ ditulis $\mathrm{char}(R)$. Gelanggang lawan dari $R$
ditulis $R^{\opp}$ dan mempunyai urutan perkalian yang berlawanan dengan $R$.
% Koreksi O014-C002: sumber mengindeks A^{\opp}, padahal paragraf ini
% mendefinisikan R^{\opp}.
\index[sym1]{Rop@$R^{\opp}$}

% segment-id: o014.li2.prelude.conventions.p035
Misalkan $r$ unsur tak nol dari gelanggang komutatif $R$. Jika homomorfisme grup
aditif $R \xrightarrow{\text{perkalian dengan}\; r} R$ mempunyai kernel tak
nol, maka $r$ disebut pembagi nol. Ideal dalam gelanggang komutatif yang
dihasilkan oleh unsur-unsur $x, y, \ldots$ ditulis $(x, y, \ldots)$.

% segment-id: o014.li2.prelude.conventions.p036
Aljabar polinomial atas gelanggang komutatif $R$ dengan peubah
$X_1, X_2, \ldots$ ditulis $R[X_1, X_2, \ldots]$, sedangkan gelanggang deret
pangkat formal ditulis $R\llbracket X_1, X_2, \ldots \rrbracket$. Jika $M$
suatu monoid, aljabar monoid atas $R$ yang terkait ditulis $R[M]$.

% segment-id: o014.li2.prelude.conventions.p037
Secara konvensi, $\Z \subset \Q \subset \R \subset \CC$ berturut-turut
menyatakan gelanggang bilangan bulat, lapangan bilangan rasional, lapangan
bilangan real, dan lapangan bilangan kompleks. Jika $q$ merupakan pangkat suatu
bilangan prima, maka $\F_q$ menyatakan lapangan hingga dengan $q$ unsur. Jika
$F \hookrightarrow E$ suatu pembenaman lapangan, perluasan lapangan terkait,
atau lapangan perluasannya, ditulis $E|F$ dan derajatnya ditulis $[E:F]$.
Jika $E|F$ merupakan perluasan
Galois, grup Galoisnya ditulis $\Gal(E|F)$.
\index[sym1]{Gal}

% Blok matriks berikut tidak aktif dalam sumber dan tetap dinonaktifkan.
%Buku ini memakai konvensi baris mendatar dan kolom tegak; matriks $n \times m$ ditulis
%\[ A = (a_{ij})_{\substack{1 \leq i \leq n \\ 1 \leq j \leq m}} = \begin{tikzpicture}[baseline]
%	\matrix (M) [matrix of math nodes, left delimiter=(, right delimiter=)] {
%		& \vdots & \\
%		\cdots & a_{ij} & \cdots \\
%		& \vdots & \\
%	};
%	\node[right=2.5em] at (M-2-3) {\scriptsize baris ke-$i$};
%	\node[below=1em] at (M-3-2) {\scriptsize kolom ke-$j$};
%\end{tikzpicture}\]
%transposenya ditulis ${}^t A$; matriks identitas $n \times n$ ditulis $1_{n \times n}$ atau $1$. Untuk matriks $n \times n$ atas gelanggang komutatif, determinan $A$ ditulis $\det A$.

% segment-id: o014.li2.prelude.conventions.p038
Notasi bagi kategori-kategori yang umum dari beberapa struktur aljabar dasar
adalah sebagai berikut.

% segment-id: o014.li2.prelude.conventions.t001
\begin{center}\small \begin{tabular}{|c|c|}\hline
	Monoid & $\cate{Mon}$ \\
	Grup &  $\cate{Grp}$ \\
	Grup abelian & $\cate{Ab} = \Z\dcate{Mod}$ \\
	Modul kiri atau kanan atas $R$ & $R\dcate{Mod}$ atau $\cated{Mod}R$ \\
	Ruang vektor atas $\Bbbk$ & $\cate{Vect}(\Bbbk)$ \\
	Bimodul $(R,S)$ & $(R,S)\dcate{Mod}$ \\
	Aljabar atas $\Bbbk$ & $\Bbbk\dcate{Alg}$ \\
	Aljabar komutatif atas $\Bbbk$ & $\Bbbk\dcate{CAlg}$
	\\ \hline
\end{tabular} \quad \begin{tabular}{l}
	$R$, $S$: gelanggang sebarang, \\
	$\Bbbk$: lapangan (untuk ruang vektor) \\
	$\Bbbk$: gelanggang komutatif (untuk aljabar) \\
	kategori modul $\Bbbk\dcate{Mod}$ tidak membedakan kiri dan kanan.
\end{tabular}\end{center}
\index[sym1]{Mod@$\cate{Mod}$}
\index[sym1]{Ab@$\cate{Ab}$}
\index[sym1]{Grp@$\cate{Grp}$}
\index[sym1]{Vect@$\cate{Vect}$}

% segment-id: o014.li2.prelude.conventions.p039
Jika $\Bbbk$ telah dipilih dan $R$ serta $S$ merupakan aljabar atas $\Bbbk$,
maka ketika membahas bimodul $(R,S)$ secara baku struktur bimodul dianggap
berasal dari struktur modul kiri atas $R \dotimes{\Bbbk} S^{\opp}$. Dengan kata lain,
perkalian kiri dan kanan oleh $\Bbbk$ disyaratkan berimpit.

% segment-id: o014.li2.prelude.conventions.p040
Grup homomorfisme antarsesama modul kiri, ataupun antarsesama modul kanan, atas
$R$ ditulis dalam bentuk $\Hom_R(X, Y)$. Untuk bimodul $(R,S)$ digunakan notasi serupa,
$\Hom_{(R, S)}(X, Y)$.

% segment-id: o014.li2.prelude.conventions.p041
Untuk gelanggang komutatif $R$ dan himpunan multiplikatifnya $U$, lokalisasi
terkait ditulis $R[U^{-1}]$.

% segment-id: o014.li2.prelude.conventions.p042
Seperti kategori himpunan $\cate{Set}$, grup, gelanggang, modul, dan struktur
lain di sini secara baku direalisasikan pada himpunan-$\mathcal{U}$, dengan
$\mathcal{U}$ semesta Grothendieck yang telah dipilih. Karena itu, setiap
kategori dalam tabel mempunyai funktor pelupa menuju $\cate{Set}$. Contoh lain
yang dibahas dalam buku ini, seperti kategori ruang topologis $\cate{Top}$ dan
kategori ruang Hausdorff kompak $\cate{CHaus}$, tunduk pada asumsi yang sama:
ruang-ruang yang dibahas direalisasikan pada himpunan-$\mathcal{U}$.

% segment-id: o014.li2.prelude.conventions.h008
\subsection*{Adjungsi}

% segment-id: o014.li2.prelude.conventions.p043
Pasangan adjoin $(F, G, \varphi)$ terdiri atas sepasang funktor berikut.

% segment-id: o014.li2.prelude.conventions.d004
$\begin{tikzcd}
	\mathcal{C} \arrow[shift left, r, "F"] & \mathcal{C}' \arrow[shift left, l, "G"]
\end{tikzcd}$

% segment-id: o014.li2.prelude.conventions.p044
Pasangan tersebut juga dilengkapi dengan suatu keluarga bijeksi kanonik
$\varphi_{X, Y}: \Hom_{\mathcal{C}'}\left(FX, Y \right) \rightiso
\Hom_{\mathcal{C}}\left(X, GY \right)$, yakni morfisme antarfunktor berikut:

% segment-id: o014.li2.prelude.conventions.d005
\[ \varphi: \Hom_{\mathcal{C}'}\left(F(\cdot), \cdot \right) \rightiso
\Hom_{\mathcal{C}}\left(\cdot, G(\cdot) \right): \;
\mathcal{C}^{\opp} \times \mathcal{C}' \to \cate{Set}. \]

% segment-id: o014.li2.prelude.conventions.p045
Pasangan adjoin juga sering ditulis dengan notasi berikut.

% segment-id: o014.li2.prelude.conventions.d006
$\begin{tikzcd}
	F: \mathcal{C} \arrow[shift left, r] & \mathcal{C}' \arrow[l, shift left] : G
\end{tikzcd}$

% segment-id: o014.li2.prelude.conventions.p046
Notasi ini menyatakan bahwa $F$ adalah adjoin kiri dari $G$, atau bahwa $G$ adalah
adjoin kanan dari $F$. Data $\varphi$ sering dihilangkan, dan pasangan adjoin
itu kemudian ditulis $(F, G)$.

% segment-id: o014.li2.prelude.conventions.p047
Dalam pasangan adjoin, $\varphi$ secara ekuivalen dapat dicirikan melalui
morfisme \emph{unit} $\eta: \identity_{\mathcal{C}} \to GF$ dan morfisme
\emph{kounit} $\varepsilon: FG \to \identity_{\mathcal{C}'}$. Agar
$(F, G, \eta, \varepsilon)$ memberikan pasangan adjoin, syarat perlu dan
cukupnya ialah bahwa morfisme-morfisme tersebut memenuhi
\emph{identitas segitiga}
\index{unit (unit)}
\index{kounit (counit)}
\index{identitas segitiga (triangle identities)}:

% segment-id: o014.li2.prelude.conventions.d007
\[ G\varepsilon \circ \eta G = \identity_G, \quad
\varepsilon F \circ F\eta = \identity_F. \]

% segment-id: o014.li2.prelude.conventions.p048
Lihat \cite[(2.6) + Proposisi~2.6.5]{Li1}. Secara ekuivalen, identitas tersebut
dapat ditulis sebagai persamaan komposisi $2$-sel berikut:

% segment-id: o014.li2.prelude.conventions.d008
\begin{equation*}\begin{gathered}
	\begin{tikzcd}
		\mathcal{C}' \arrow[r, "G"] \arrow[rr, bend right=60, "\identity_{\mathcal{C}'}"', "" {name=B}] & \mathcal{C} \arrow[r, "F"] \arrow[Rightarrow, to=B, "\varepsilon"] \arrow[rr, bend left=60, "\identity_{\mathcal{C}}", ""' {name=T}] & \mathcal{C}' \arrow[r, "G"] \arrow[Rightarrow, from=T, "\eta"] & \mathcal{C}
	\end{tikzcd} \quad = \quad \begin{tikzcd}[column sep=large]
		\mathcal{C}' \arrow[r, bend left=60, "G", ""' {name=T}] \arrow[r, bend right=60, "G"', "" {name=B}] & \mathcal{C} \arrow[Rightarrow, from=T, to=B, "{\identity_G}" description]
	\end{tikzcd} , \\
	\begin{tikzcd}
		\mathcal{C} \arrow[r, "F"] \arrow[rr, bend left=60, "\identity_{\mathcal{C}}", ""' {name=T}] & \mathcal{C}' \arrow[r, "G"] \arrow[Rightarrow, from=T, "\eta"] \arrow[rr, bend right=60, "\identity_{\mathcal{C}'}"', ""' {name=B}] & \mathcal{C} \arrow[r, "F"] \arrow[Rightarrow, to=B, "\varepsilon"] & \mathcal{C}'
	\end{tikzcd} \quad = \quad \begin{tikzcd}[column sep=large]
		\mathcal{C} \arrow[r, bend left=60, "F", ""' {name=T}] \arrow[r, bend right=60, "F"', "" {name=B}] & \mathcal{C}' \arrow[Rightarrow, from=T, to=B, "{\identity_F}" description]
	\end{tikzcd}.
\end{gathered}\end{equation*}

% segment-id: o014.li2.prelude.conventions.p049
Untuk rinciannya, lihat \cite[Catatan~2.6.6]{Li1}.

% segment-id: o014.li2.prelude.conventions.h009
\subsection*{Limit}

% segment-id: o014.li2.prelude.conventions.p050
Perhatikan funktor $\alpha: I \to \mathcal{C}$.

% segment-id: o014.li2.prelude.conventions.l005
\begin{itemize}
	% segment-id: o014.li2.prelude.conventions.i015
	\item Limit induktif $\alpha$, jika ada, ditulis $\varinjlim \alpha$ atau
	$\varinjlim_{i \in \Obj(I)} \alpha(i)$. Limit ini merupakan objek
	$\mathcal{C}$ yang dilengkapi suatu keluarga morfisme
	$\alpha(i) \to \varinjlim \alpha$ dan ditentukan oleh sifat universal.
	\index{limit (limit)}
	% segment-id: o014.li2.prelude.conventions.i016
	\item Demikian pula, limit projektif $\alpha$, jika ada, ditulis
	$\varprojlim \alpha$ atau $\varprojlim_{i \in \Obj(I)} \alpha(i)$. Limit
	ini dilengkapi suatu keluarga morfisme $\varprojlim \alpha \to \alpha(i)$
	dan ditentukan oleh sifat universal.
\end{itemize}

% segment-id: o014.li2.prelude.conventions.p051
Kedua jenis limit tersebut saling dual: $\varinjlim$ dari
$\alpha: I \to \mathcal{C}$ dan $\varprojlim$ dari
$\alpha^{\opp}: I^{\opp} \to \mathcal{C}^{\opp}$ merupakan hal yang sama.
Oleh karena itu, dalam buku ini funktor yang terkait dengan $\varprojlim$ juga
sering ditulis dalam bentuk $I^{\opp} \to \mathcal{C}$.

% segment-id: o014.li2.prelude.conventions.p052
Dalam banyak pustaka, $\varinjlim$ juga disebut kolimit dan ditulis
$\mathrm{colim}$, sedangkan $\varprojlim$ disebut limit dan ditulis $\lim$.
Kecuali dinyatakan lain, istilah ``limit'' dalam buku ini mencakup
$\varinjlim$ maupun $\varprojlim$.
\index[sym1]{lim@$\varinjlim$, $\varprojlim$, $\lim$, $\mathrm{colim}$}

% segment-id: o014.li2.prelude.conventions.p053
Jika $I$ adalah kategori kecil, maka $\varinjlim$ atau $\varprojlim$ terkait disebut
\emph{limit kecil}\index{limit!kecil (small)}. Mengikuti terminologi standar
dalam \cite[\S~2.8]{Li1}, jika kategori $\mathcal{C}$ mempunyai semua
$\varprojlim$ kecil (atau semua $\varinjlim$ kecil), maka $\mathcal{C}$
disebut \emph{lengkap} (atau \emph{kolengkap}). Pengertian ini bergantung pada
pilihan $\mathcal{U}$.
\index{kategori!lengkap, kolengkap (complete, cocomplete)}

% segment-id: o014.li2.prelude.conventions.p054
Untuk menetapkan notasi, selanjutnya kita memilih kategori $\mathcal{C}$ dan
meninjau kembali beberapa kasus khusus limit yang lazim. Definisi terperinci
dapat ditemukan dalam \cite[\S~2.7]{Li1} atau buku ajar lain. Dalam
sifat-sifat universal berikut, $T$ menyatakan sebarang objek dalam
$\mathcal{C}$, sedangkan $I$ menyatakan sebarang himpunan.

% segment-id: o014.li2.prelude.conventions.t002
\begin{center}\small\setlength{\tabcolsep}{2.5pt}\begin{tabular}{|c|c|c|c|c|} \hline
	Nama limit & Masukan & Objek & Morfisme kanonik & Sifat universal \\ \hline
	Penyama & $\begin{tikzcd}[column sep=small] X \arrow[shift left, r, "f"] \arrow[shift right, r, "g"'] & Y\end{tikzcd}$ & $\Ker(f,g)$ & $\Ker(f,g) \xrightarrow{\iota} X$ & $\begin{tikzcd}[row sep=small, column sep=tiny, every cell/.append style={font = \footnotesize}]
		\Hom(T, \Ker(f,g)) \arrow[d, "1:1"'] & \psi \arrow[mapsto, d] \\
		\left\{\begin{array}{l} \phi \in \Hom(T, X): \\ f\phi = g\phi \end{array}\right\} & \iota\psi
	\end{tikzcd}$ \\
	Kopenyama & (seperti di atas) & $\Coker(f,g)$ & $Y \xrightarrow{p} \Coker(f,g)$ & $\begin{tikzcd}[row sep=small, column sep=tiny, every cell/.append style={font = \footnotesize}]
		\Hom(\Coker(f,g), T) \arrow[d, "1:1"'] & \psi \arrow[mapsto, d] \\
		\left\{\begin{array}{l} \phi \in \Hom(Y, T): \\ \phi f = \phi g \end{array}\right\} & \psi p
	\end{tikzcd}$ \\
	Produk & $(X_i)_{i \in I}$ & $\displaystyle\prod_{i \in I} X_i$ & $\displaystyle\prod_{j \in I} X_j \xrightarrow{p_i} X_i$ & $\begin{tikzcd}[row sep=small, column sep=tiny, every cell/.append style={font = \footnotesize}]
		\Hom\left(T, \displaystyle\prod_{i \in I} X_i \right) \arrow[d, "1:1"'] & \psi \arrow[mapsto, d] \\
		\displaystyle\prod_{i \in I} \Hom(T, X_i) & \left( p_i \psi \right)_{i \in I}
	\end{tikzcd}$ \\
	Koproduk & (seperti di atas) & $\displaystyle\coprod_{i \in I} X_i$ & $X_i \xrightarrow{\iota_i} \displaystyle\coprod_{j \in I} X_j$ & $\begin{tikzcd}[row sep=small, column sep=tiny, every cell/.append style={font = \footnotesize}]
		\Hom\left(\displaystyle\coprod_{i \in I} X_i, T \right) \arrow[d, "1:1"'] & \psi \arrow[mapsto, d] \\
		\displaystyle\prod_{i \in I} \Hom(X_i, T) & \left( \psi \iota_i \right)_{i \in I}
	\end{tikzcd}$ \\ \hline
\end{tabular}\end{center}

% segment-id: o014.li2.prelude.conventions.n001
\begin{quote}\small
\textbf{Catatan penerjemah.} Pada baris koproduk, sumber menulis sasaran
injeksi kanonik sebagai $\prod_{j \in I}X_j$. Notasi itu dikoreksi menjadi
$\coprod_{j \in I}X_j$ sesuai objek pada kolom sebelumnya dan sifat universal
di kolom berikutnya; lihat koreksi O014-C003.
\end{quote}

% segment-id: o014.li2.prelude.conventions.p055
Produk hingga (atau koproduk hingga) juga ditulis sebagai
$X_1 \times X_2 \times \cdots$ (atau $X_1 \sqcup X_2 \sqcup \cdots$). Jika
$\mathcal{C}$ kategori aditif, koproduk $\coprod_i$ sering ditulis dengan
simbol jumlah langsung $\bigoplus_i$ dari teori modul.
\index[sym1]{oplus@$\oplus$}

% segment-id: o014.li2.prelude.conventions.p056
Penyama/kopenyama dan produk/koproduk masing-masing merupakan pasangan dual.
Pencirian melalui sifat universal tersebut juga dapat ditulis sebagai diagram
komutatif seperti dalam \cite[\S~2.7]{Li1}, atau ditafsirkan melalui pembenaman
Yoneda yang akan ditinjau dalam \S\sourcecrossref{sec:Yoneda}{A.1}. Andaikan
$\coprod_{i \in I} X_i$ dan $\prod_{j \in J} Y_j$ ada; sifat universal
memberikan bijeksi

% segment-id: o014.li2.prelude.conventions.d009
\[ \Hom\left( \coprod_{i \in I} X_i, \prod_{j \in J} Y_j \right)
\xrightarrow{1:1} \prod_{\substack{i \in I \\ j \in J}} \Hom(X_i, Y_j),
\quad \psi \mapsto (p_j \psi \iota_i)_{i,j} . \]

% segment-id: o014.li2.prelude.conventions.l006
\begin{description}
	% segment-id: o014.li2.prelude.conventions.i017
	\item[Objek terminal] Produk yang terkait dengan $I = \emptyset$, jika ada,
	disebut objek terminal dalam $\mathcal{C}$ dan untuk sementara ditulis $Z$.
	Menurut konvensi, $\prod_{i \in \emptyset} := \{\emptyset\}$, sehingga
	sifat universal $Z$ adalah: $\Hom(T, Z)$ merupakan himpunan satu unsur untuk
	setiap $T \in \Obj(\mathcal{C})$.
	% segment-id: o014.li2.prelude.conventions.i018
	\item[Objek awal] Koproduk yang terkait dengan $I = \emptyset$, jika ada,
	disebut objek awal dalam $\mathcal{C}$ dan untuk sementara ditulis $S$.
	Sifat universalnya adalah: $\Hom(S, T)$ merupakan himpunan satu unsur untuk
	setiap $T \in \Obj(\mathcal{C})$.
	% segment-id: o014.li2.prelude.conventions.i019
	\item[Objek nol dan morfisme nol] Jika objek $0$ dalam $\mathcal{C}$ sekaligus
	merupakan objek awal dan objek terminal, objek itu disebut objek nol;
	morfisme ke dan dari objek nol ada dan tunggal. Untuk setiap
	$X, Y \in \Obj(\mathcal{C})$, morfisme komposit $X \to 0 \to Y$
	memberikan unsur kanonik dalam $\Hom(X, Y)$ yang disebut morfisme nol.
	\index{objek!nol (zero)}
\end{description}

% segment-id: o014.li2.prelude.conventions.p057
Karena setiap isomorfisme dari objek $X$ ke objek nol, jika ada, bersifat
tunggal, lazim dikatakan bahwa $X$ adalah objek nol atau ditulis $X = 0$,
alih-alih mengatakan bahwa $X$ isomorfik dengan objek nol.

% segment-id: o014.li2.prelude.conventions.p058
Selanjutnya kita meninjau dua jenis limit penting yang saling dual.

% segment-id: o014.li2.prelude.conventions.t003
\begin{center}\small\setlength{\tabcolsep}{2.5pt}\begin{tabular}{|c|c|c|c|c|} \hline
	Nama & Masukan & Objek & Morfisme kanonik & Sifat universal \\ \hline
	Produk serat & $\left( X_i \xrightarrow{f_i} Z \right)_{i \in I}$ & $\displaystyle\prod_{i \in I} (X_i \to Z)$ & $\begin{tikzcd}[every cell/.append style={font = \footnotesize}] \displaystyle\prod_{j \in I} (X_j \to Z) \arrow[d, "p_i"] \\ X_i \end{tikzcd}$ & $\begin{tikzcd}[every cell/.append style={font = \footnotesize}]
		\Hom\left(T, \displaystyle\prod_{i \in I} (X_i \to Z) \right) \arrow[d, "1:1"', "{\psi \mapsto (p_i \psi)_i}"] \\
		\left\{\begin{array}{l} (\xi_i: T \to X_i)_{i \in I}: \\ \forall i,j \in I, \\ f_i \xi_i = f_j \xi_j \end{array}\right\}
	\end{tikzcd}$ \\
	Koproduk serat & $\left( Z \xrightarrow{g_i} X_i \right)_{i \in I}$ & $\displaystyle\coprod_{i \in I} (Z \to X_i)$ & $\begin{tikzcd}[every cell/.append style={font = \footnotesize}] X_i \arrow[d, "\iota_i"] \\ \displaystyle\coprod_{j \in I} (Z \to X_j) \end{tikzcd}$ & $\begin{tikzcd}[every cell/.append style={font = \footnotesize}]
		\Hom\left(\displaystyle\coprod_{i \in I} (Z \to X_i), T \right) \arrow[d, "1:1"', "{\psi \mapsto (\psi \iota_i)_i}"] \\
		\left\{\begin{array}{l} (\xi_i: X_i \to T)_{i \in I}: \\ \forall i,j \in I, \\ \xi_i g_i = \xi_j g_j \end{array}\right\}
	\end{tikzcd}$ \\ \hline
\end{tabular}\end{center}

% segment-id: o014.li2.prelude.conventions.p059
Dengan mengambil $\psi = \identity$ dalam sifat universal produk serat,
terlihat bahwa $f_i p_i = f_j p_j$ untuk semua $i,j$; dari sini diperoleh
morfisme kanonik $\prod_{i \in I} (X_i \to Z) \to Z$. Demikian pula, dengan
mengambil $\psi = \identity$ untuk koproduk serat, terlihat bahwa
$\iota_i g_i = \iota_j g_j$ untuk semua $i,j$; dari sini diperoleh morfisme
kanonik $Z \to \coprod_{i \in I} (Z \to X_i)$.

% segment-id: o014.li2.prelude.conventions.p060
Produk serat yang melibatkan sejumlah berhingga objek
$X_1 \dtimes{Z} X_2 \dtimes{Z} \cdots$, maupun koproduk serat
$X_1 \dsqcup{Z} X_2 \dsqcup{Z} \cdots$, dapat direduksi ke kasus dua objek.
Perhatikan dua diagram komutatif berikut.

% segment-id: o014.li2.prelude.conventions.d010
\[\begin{tikzcd}
	X \dtimes{Z} Y \arrow[r, "p_1"] \arrow[d, "p_2"'] & X \arrow[d] \\
	Y \arrow[r] & Z
\end{tikzcd} \quad \text{dan} \quad \begin{tikzcd}
	Z \arrow[r] \arrow[d] & X \arrow[d, "\iota_1"] \\
	Y \arrow[r, "\iota_2"'] & X \dsqcup{Z} Y
\end{tikzcd}\]
% Koreksi O014-C004: garis miring terbalik tambahan setelah pembuka tikzcd
% pada sumber dihapus sebagai artefak sintaksis tanpa muatan matematika.

% segment-id: o014.li2.prelude.conventions.p061
Kedua diagram tersebut masing-masing disebut diagram \emph{tarik balik} dan diagram
\emph{dorong keluar}. Morfisme $X \dtimes{Z} Y \to Y$ disebut tarik balik dari
$X \to Z$, sedangkan $Y \to X \dsqcup{Z} Y$ disebut dorong keluar dari
$Z \to X$; perhatikan bahwa peran $X$ dan $Y$ simetris. Lebih umum, diagram
komutatif yang isomorfik dengan diagram-diagram di atas juga masing-masing
disebut diagram tarik balik dan diagram dorong keluar; lihat
\cite[Definisi~2.8.5]{Li1}.
\index{diagram tarik balik (pullback diagram)}
\index{diagram dorong keluar (pushout diagram)}

% segment-id: o014.li2.prelude.conventions.p062
Dalam notasi buku ini, pusat persegi pada diagram tarik balik (atau dorong
keluar) diberi tanda $\Box$ (atau $\boxplus$).
\index[sym1]{Box@$\Box$, $\boxplus$}

% segment-id: o014.li2.prelude.conventions.h010
\subsection*{Penerjemahan}

% segment-id: o014.li2.prelude.conventions.p063
Dalam edisi sumber berbahasa Tionghoa, konvensi penerjemahan sama dengan
\cite{Li1}; pilihan istilah berbahasa Tionghoa tetap diupayakan selaras dengan
\cite{ZG}, kecuali jika terjemahan dalam rujukan tersebut jelas kurang tepat
atau keliru.

% segment-id: o014.li2.prelude.conventions.p064
Dalam indeks sumber, istilah berbahasa Tionghoa disertai padanan bahasa
Inggris. Pertama, hal ini membantu pembaca yang perlu membaca atau menulis
artikel berbahasa Inggris; kedua, memberikan rujukan bagi pembaca yang bahasa
ibunya bukan bahasa Tionghoa; terakhir, karena alasan historis, penguasaan
bahasa asing masih diperlukan untuk memahami makna banyak simbol matematika.

% segment-id: o014.li2.prelude.conventions.n002
\begin{quote}\small
\textbf{Catatan penerjemah.} Edisi Indonesia ini mempertahankan padanan bahasa
Inggris dalam indeks dan mencatat pilihan istilah Indonesia--Inggris dalam
glosarium produksi yang terpisah.
\end{quote}
